\scanpage{10}

\figure{\begin{tikzpicture}\node[draw,ellipse,inner sep=3pt]{2};\end{tikzpicture}}
Supposons que $U$ soit \deleted{muni d'une} \inserted{une} \figure{\begin{tikzpicture}\node[draw,ellipse,inner sep=2pt] at (0,0){$k$};\node at (0,-.6){$k$};\end{tikzpicture}}-algèbre
augmentée ($k$ anneau comm.) avec un
hom diagonal
\[
\Delta:U\;\text{---}\;U\widehat\otimes_k U
\]
qui soit un hom de $k$-algèbres augmentées.
$U$ est appelée une $k$-hyperalgèbre \inserted{topologique} (au sens
large). On désigne par
\deleted{$\unsure{g(U)}$} \inserted{$\unsure{\gamma(U)|G(U)}$}
\deleted{$\unsure{G(U)}$} l'ens
des $g\in1+I(U)$ \deleted{telles que} \inserted{telles} que
\[
\Delta(g)=g\otimes g
\]
i.e.
\[
\delta\xi=\xi\otimes\xi\qquad\text{si }g=1+\xi\quad(\xi\in I(U))
\]
\margin{on suppose que $I(U)^n\longrightarrow0$ pour $n\to\infty$.
\deleted{\unsure{Donc}}}
[\deleted{éléments} ils correspondent par dualité aux
hom de la hyperalgèbre \deleted{\unsure{triviale}}
$k[T]$, avec $\delta T=T\otimes T$ ; on peut aussi
prendre $k[[T]]$, avec $\delta T=T\otimes T$, et
se borner aux hom. continus\ldots ;
cette \inserted{hyper}Algèbre est le dual top. de
l'Algèbre du groupe $\mathbb Z$ ..]. \deleted{\unsure{Ils}} Les
$g$ forment un groupe, car
\[
\Delta(gg')=\Delta(g)\Delta(g')=(g\otimes g)(g'\otimes g')
=gg'\otimes\unsure{gg'}
\]
\[
\Delta(g^{-1})=\deleted{\unsure{g\otimes}}\Delta(g)^{-1}
=(g\otimes g)^{-1}=g^{-1}\otimes g^{-1}.
\]
Fonctorialité de $\unsure{\gamma(U)|G(U)}$, et $\unsure{\pi(U)|P(U)}$
(voir ci-dessous)
\deleted{\unsure{par}} \deleted{\unsure{les}}
$(k,U)\longrightarrow(k',U')$.

\mspage{9}
